% ═══════════════════════════════════════════════════════════════════════════════
%  CAPÍTULO — Visitor
% ═══════════════════════════════════════════════════════════════════════════════

\chapter{Visitor}

% ─── Situación Problema ──────────────────────────────────────────────────────
\section{Situación Problema}

El equipo de analítica de Airbnb necesita ejecutar distintos tipos de
análisis sobre el catálogo de alojamientos: calcular el ingreso
potencial máximo de cada propiedad, evaluar el índice de cumplimiento
de estándares de calidad y generar un reporte fiscal con la información
relevante para declaraciones tributarias en cada país. Cada análisis
opera sobre los mismos objetos del dominio —\texttt{Habitacion},
\texttt{Apartamento}, \texttt{CasaRural}— pero extrae y procesa datos
distintos.

Añadir la lógica de cada análisis directamente a las clases del dominio
las contamina con responsabilidades que no le corresponden. Cada vez que
el equipo de analítica requiere un nuevo tipo de informe, hay que modificar
todas las clases de alojamiento, rompiendo el principio de
responsabilidad única y arriesgando errores en el núcleo del modelo de
dominio.

% ─── Solución ────────────────────────────────────────────────────────────────
\section{Solución}

El patrón Visitor separa los algoritmos de análisis de las clases sobre
las que operan. Cada clase de alojamiento implementa un método
\texttt{aceptar(visitante)} que simplemente invoca el método correspondiente
del visitante. Cada análisis se implementa como un visitante concreto
—\texttt{VisitanteIngresoMaximo}, \texttt{VisitanteCalidad},
\texttt{VisitanteFiscal}— con métodos específicos para cada tipo de
alojamiento.

Las clases del dominio no necesitan cambiar cuando se requiere un nuevo
análisis; basta con crear un nuevo visitante. Cada visitante encapsula
de forma cohesiva toda la lógica de su análisis particular, y la
estructura del modelo de dominio permanece limpia y estable.

% ─── Diagrama de clases ─────────────────────────────────────────────────────
\begin{figure}[H]
    \centering
    % \includegraphics[width=\textwidth]{imagenes/abstract_factory_diagrama.png}
    \caption{Diagrama de clases del patrón Abstract Factory}
\end{figure}


% ─── Roles de las Clases ────────────────────────────────────────────────────
\section{Roles de las Clases}

% Explicar la responsabilidad de cada clase/interfaz

\begin{itemize}
    \item \textbf{AbstractFactory:}
    \item \textbf{ConcreteFactory:}
    \item \textbf{AbstractProduct:}
    \item \textbf{ConcreteProduct:}
    \item \textbf{Client:}
\end{itemize}


% ─── Main ───────────────────────────────────────────────────────────────────
\section{Main}

% Explicar qué realiza el programa principal

\begin{lstlisting}[language=Java, caption=NombreClase]
 
\end{lstlisting}


% ─── Salida del Main ────────────────────────────────────────────────────────
\section{Salida del Main}

\begin{lstlisting}[language=Java, caption=NombreClase]
 
\end{lstlisting}


% ─── Clases ─────────────────────────────────────────────────────────────────
\section{Clases}

% ─── Clase 1 ────────────────────────────────────────────────────────────────
\subsection{NombreClase}

\begin{lstlisting}[language=Java, caption=NombreClase]
 
\end{lstlisting}


% ─── Clase 2 ────────────────────────────────────────────────────────────────
\subsection{NombreClase}

\begin{lstlisting}[language=Java, caption=NombreClase]
 
\end{lstlisting}


% ─── Clase 3 ────────────────────────────────────────────────────────────────
\subsection{NombreClase}

\begin{lstlisting}[language=Java, caption=NombreClase]
 
\end{lstlisting}


% ─── Conclusiones ───────────────────────────────────────────────────────────
\section{Conclusiones}

% Explicar:
% - Ventajas del patrón
% - Cuándo usarlo
% - Beneficios obtenidos
% - Posibles desventajas
% - Relación con principios SOLID